\subsection{Nonlinear Variational Problems}\label{sec:nonlinear-variational-problems}

\subsubsection{Convex problems}\label{subsec:convex-variational-problems}

An important class of nonlinear elliptic PDEs arise from problems in the calculus of variations. We will focus on a particularly vital set of representatives of this class: problems describing the equilibrium of \emph{elastic} materials. In discussing this class of problems we will be able to examine not only many of the classical techniques for elliptic problems from the calculus of variations, but also some important techniques which have been developed over the last twenty years.

We begin our brief discussion of nonlinear elasticity by describing the \emph{kinematics} or geometry of deformation of three-dimensional bodies. We let a domain $\Omega\subset\mathbb{R}^3$ represent the \emph{reference configuration} of a material body. We assume that $\Omega$ is bounded with Lipschitz boundary. (It is usually convenient to think of the reference configuration $\Omega$ as the ``rest\textquotesingle\textquotesingle\ or ``unstressed\textquotesingle\textquotesingle\ configuration of the body, but we will not restrict ourselves to this case.) A \emph{deformation} of the body is simply a mapping of the form
\begin{equation}\label{eq:elasticity-reference-deformation-mapping}
\mathbb{R}^3\supset\Omega\ni\bm{x}\mapsto\bm{p}(\bm{x})\in\mathbb{R}^3.
\end{equation}
It will ease our computations somewhat to spell out all of our vectors and matrices in terms of components. Thus, we let $\bm{e}_1,\bm{e}_2,\bm{e}_3$ constitute a fixed orthonormal basis for $\mathbb{R}^3$, and for a given vector $\bm{v}$, we define $v_i:=\bm{v}\cdot\bm{e}_i$.

We will assume a certain amount of smoothness of the deformation: We either make the classical assumption that $\bm{p}\in C_b^1(\overline\Omega)$ or use a Sobolev space $\bm{p}\in W^{1,p}(\Omega)$. In either case, we are able to define the \emph{deformation gradient}
\begin{equation}\label{eq:elasticity-deformation-gradient}
\bm{F}(\bm{x}):=\dfrac{\partial\bm{p}}{\partial\bm{x}}(\bm{x}).
\end{equation}
In terms of components, we have
\begin{equation}\label{eq:elasticity-deformation-gradient-components}
F_{ij}(\bm{x}):=\dfrac{\partial p_i}{\partial x_j}(\bm{x}).
\end{equation}
We would like to restrict the deformations to be one-to-one mappings so that the material does not interpenetrate or overlap when it deforms. However, such a constraint is hard to treat analytically so we will ignore it. We will, however, assume that the material \emph{preserves orientation} (i.e., no ``mirror-image\textquotesingle\textquotesingle\ deformations). If $\bm{p}\in C_b^1(\overline\Omega)$ this constraint can be expressed by the pointwise inequality
\begin{equation}\label{eq:elasticity-positive-orientation-determinant}
\det\bm{F}(\bm{x})>0.
\end{equation}
If we are taking $\bm{p}\in W^{1,p}(\Omega)$, then we will assume that the inequality \eqref{eq:elasticity-positive-orientation-determinant} is satisfied almost everywhere.

We will specify \emph{displacement boundary conditions} on our deformations. The most natural way to state such conditions is to specify a continuous function $\bm{b}:\partial\Omega\to\mathbb{R}^3$ and to require that
\begin{equation}\label{eq:elasticity-prescribed-boundary-deformation}
\bm{p}(\bm{x})=\bm{b}(\bm{x})\quad\text{for all }\bm{x}\in\partial\Omega.
\end{equation}
If we wish to consider $\bm{p}\in W^{1,p}(\Omega)$ we have to require that \eqref{eq:elasticity-prescribed-boundary-deformation} holds in the sense of trace. However, when working in Sobolev spaces, it is more convenient to enforce boundary conditions by specifying a function $\bm{g}:\Omega\to\mathbb{R}^3$ with $\bm{g}\in W^{1,p}(\Omega)$ and requiring
\begin{equation}\label{eq:elasticity-homogeneous-trace-difference}
\bm{p}_0:=\bm{p}-\bm{g}\in W_0^{1,p}(\Omega).
\end{equation}
This ensures that $\bm{p}$ and $\bm{g}$ have the same trace on the boundary of $\Omega$. If we are working with deformations in the space $H^1(\Omega)$ and $\partial\Omega$ is of class $C^1$, then Theorem~\ref{thm:higher-normal-trace-bounded-right-inverse} (the inverse trace theorem) ensures that if $\bm{b}\in H^{1/2}(\partial\Omega)$, then there exists $\bm{g}\in H^1(\Omega)$ such that the trace of $\bm{g}$ on $\partial\Omega$ is $\bm{b}$. Thus, by using this $\bm{g}$, the two boundary conditions \eqref{eq:elasticity-prescribed-boundary-deformation} and \eqref{eq:elasticity-homogeneous-trace-difference} can be made to coincide. We can achieve a similar result for the space $W^{1,p}(\Omega)$ with $p\neq2$ by using a more general version of Theorem~\ref{thm:higher-normal-trace-bounded-right-inverse}, but instead of trying to find the most general inverse trace theorem (on the roughest possible boundary) we will content ourselves with boundary conditions of the form \eqref{eq:elasticity-homogeneous-trace-difference}. In addition, we will be able to assume such conditions as
\begin{equation}\label{eq:elasticity-boundary-extension-positive-determinant}
\det\left(\dfrac{\partial\bm{g}}{\partial\bm{x}}\right)>0
\end{equation}
almost everywhere. (The existence of such a $\bm{g}$ is difficult to address in a general inverse trace theorem.) To sum up, we assume that $\bm{g}\in W^{1,p}(\Omega)$ satisfying \eqref{eq:elasticity-boundary-extension-positive-determinant} exists and take the domain of our elasticity boundary-value problem to be
\begin{equation}\label{eq:elasticity-admissible-deformation-domain}
\mathcal{D}_{\mathcal{E}}:=\{\bm{p}\in W^{1,p}(\Omega)\mid\det\nabla\bm{p}>0\text{ a.e., and }\bm{p}-\bm{g}\in W_0^{1,p}(\Omega)\}.
\end{equation}

